\clearpage
\section{JUSTIFICATIVA}
\label{JUSTIFICATIVA}

Até o século XIX, o conhecimento humano se caracterizou por uma dinâmica basicamente cumulativa, com correções ocasionais. Já no século XX, o acelerado crescimento da ciência e da tecnologia revelou uma dinâmica diferente, onde o conhecimento tecnológico passou a transformar-se velozmente. A aceleração do desenvolvimento tecnológico, aliado aos modernos processos de produção industrial, é um fenômeno que vem se difundindo mundialmente através dos processos de internacionalização e globalização das economias.

Desta forma, faz-se crescente a importância das atividades que envolvem controle e automação de processos no âmbito industrial brasileiro. Esta área atualmente vem recebendo maiores investimentos nas empresas, com o objetivo de proporcionar subsídios para uma melhor adaptação à evolução tecnológica que se impõe no novo cenário da economia mundial e passa a assumir um papel estratégico no desenvolvimento industrial, estando diretamente relacionado com a produtividade das empresas.

De acordo com o Instituto Brasileiro de Geografia e Estatística (IBGE), a taxa média de desemprego no Brasil aumentou 8,1\% no semestre encerrado em maio de 2017 e o índice da população desempregada, no período equivalente, também cresceu 18,04\%, frente ao ano passado. No contexto regional, a região Nordeste e mais especificamente o estado do Ceará, vêm experimentando um desenvolvimento expressivo com relação ao restante do país, promovendo uma expansão do mercado consumidor regional através da instalação de várias empresas nos últimos anos nas áreas de siderurgia, energia e indústria petroquímica ao redor do Complexo Portuário do Pecém e no Pólo Industrial de Maracanaú. Estas empresas precisam de profissionais com excelente formação para atuar nos cargos de liderança e engenharia. Esse fato age como importante agente fomentador do mercado produtor local que demanda cada vez mais profissionais especializados detentores de uma formação diferenciada e em consonância com a inovação tecnológica presente nos modernos processos fabris.

Segundo o Conselho Federal de Engenharia, Arquitetura e Agronomia - CONFEA, enquanto o Brasil forma cerca de 40 mil engenheiros por ano, a Rússia, a India e a China formam 190 mil, 220 mil e 650 mil, respectivamente. Entidades empresariais, como a Confederação Nacional da Indústria - CNI, têm feito estudos sobre o impacto da falta de engenheiros no desenvolvimento econômico brasileiro. E órgãos governamentais, como a Financiadora de Estudos e Projetos - FINEP, patrocinam desde 2006 programas de estímulo à formação de mais engenheiros no país. Segundo estimativas do CONFEA, o Brasil tem um déficit de 20 mil engenheiros por ano.

No país há 600 mil engenheiros, o equivalente a 6 profissionais para cada mil trabalhadores. Nos Estados Unidos e no Japão, a proporção é de 25 engenheiros por mil trabalhadores, segundo publicações da FINEP. Elas também informam que, dos 40 mil engenheiros que se diplomam anualmente no Brasil, mais da metade opta pelo curso de engenharia civil. Assim, setores como os de petróleo, gás, energia e biocombustível são os que mais sofrem com a escassez de profissionais. Além disso, os cursos de engenharia apresentam elevada taxa de evasão que em algumas instituições chega a 55\%, decorrente da elevada complexidade dos cursos e pela falta de interesse dos jovens pela profissão, ocasionada pela falta de preparo no ensino básico e médio, principalmente nas disciplinas de matemática, física e química.

Múltiplos são os indicadores que apontam para um contexto caracterizado pelo avanço tecnológico e a consequente necessidade das empresas de reverem seus processos de trabalho e, sobretudo, buscarem diferencial competitivo através de ações proativas, saindo à frente da concorrência e das expectativas do mercado. A modernização e a inovação das técnicas utilizadas, juntamente com a categoria gerencial passam, portanto, a personificar, no dia a dia empresarial, um dos diferenciais competitivos e consequentemente, de sobrevivência.

Esse contexto prima-se, portanto, pela necessidade de um profissional que atue como gerente de fábrica, empreendedor, convergindo suas atribuições técnicas específicas às atribuições de gestor; altamente qualificado com habilidades diferentes das tradicionais, preocupado em organizar tática e estrategicamente as metas a serem alcançadas pela filosofia da empresa. Um profissional apoiado na ciência e na tecnologia, motivado e motivador, e que objetive melhorias contínuas dos resultados atingidos nos processos produtivos.

Nesse contexto, o governo federal através da Lei Nº 11.892 de 29 de dezembro de 2008 institui a Rede Federal de Educação Profissional, Científica e Tecnológica e cria os Institutos Federais de Educação, Ciência e Tecnologia – Ifs. Um dos objetivos dos Institutos Federais, conforme alínea c, inciso VI, do art. 7º, é ofertar cursos em nível de educação superior, dentre eles, os cursos de bacharelado e engenharia, visando à formação de profissionais para os diferentes setores da economia e áreas do conhecimento. Portano, a Rede Federal de Ensino assume a missão de ofertar cursos de engenharia em suas unidades, como pode ser verificado no documento intitulado “Princípios Norteadores das Engenharias nos Institutos Federais”, publicado pela SETEC/MEC em abril de 2009.

A decisão em ofertar cursos de engenharia nos Institutos Federais prende-se a alguns aspectos estratégicos, considerando-se o momento singular por que passa o país e as possibilidades que a Rede Federal apresenta. Em primeiro lugar, há hoje na Rede um corpo docente com a qualificação capaz de responder ao desafio de promover a oferta desses cursos e expandir as atividades para a pesquisa, extensão e a pós-graduação. Em segundo lugar, já decorre tempo suficiente de oferta de cursos superiores nos centros federais de educação tecnológica (CEFET), para se fazer uma avaliação acerca dessa experiência e reunir elementos para os próximos desafios. Em terceiro lugar, pela oportunidade que têm os Institutos Federais de revisitar o ensino de engenharia, dentro de uma visão mais humanística e sustentável. E por fim, com vistas a atender à demanda por novos(as) engenheiros(as) oriunda das novas demandas sociais do mercado de trabalho, tendo em vista a recente retomada do desenvolvimento econômico verificado no Brasil que, em sua persistência, obrigará a um redimensionamento do setor educacional e, em particular, dos cursos de engenharia.

Atendendo a esses princípios, o IFCE, ciente dessa relevância no cenário de transformações no mundo do trabalho e na formação do cidadão e visando sua total inserção social, política, cultural e ética, tem buscado desempenhar tal tarefa com qualidade, reinterpretando o seu relacionamento com o segmento produtivo e buscando novos modelos curriculares.

Nesse contexto, o IFCE Campus Maracanaú, vem através deste projeto propor o \NOMECURSO, com vistas a formar o \NOMEPROFISSIONAL para o exercício crítico e competente da sua profissão, onde os valores e princípios estéticos, políticos e éticos sejam seus norteadores, e o estímulo à pesquisa e a postura de permanente busca de atualização profissional seja uma constante. Buscando, desta forma, assim nos termos Lei No 11.892/2008, contribuir com os diversos setores da economia, com ênfase no desenvolvimento socioeconômico local, regional e nacional.
